\section*{Capítulo I, Problema XVIII}

\textbf{Problema:}

18. Sean $B_1, \dots, B_m$ vectores de norma $1$ en el espacio $n$-dimensional, y mutuamente perpendiculares, es decir, $B_i \cdot B_j = 0$ si $i \neq j$. Sea $A$ un vector en el espacio $n$-dimensional, y sea $c_i$ la componente de $A$ a lo largo de $B_i$. Sean $x_1, \dots, x_m$ números. Demuestra que:
\[
    \|A - (c_1B_1 + \cdots + c_mB_m)\| \leq \|A - (x_1B_1 + \cdots + x_mB_m)\|.
\]

\textbf{Solución:}

Dado que $B_1, \dots, B_m$ son vectores ortogonales de norma $1$, podemos escribir $A$ como:
\[
    A = c_1B_1 + c_2B_2 + \cdots + c_mB_m + R,
\]
donde $R$ es el residuo ortogonal a todos los $B_i$, es decir, $R \cdot B_i = 0$ para todo $i$.

Sea $x = x_1B_1 + x_2B_2 + \cdots + x_mB_m$. Entonces:
\begin{align*}
    \|A - x\|^2 &= \|(c_1B_1 + \cdots + c_mB_m + R) - (x_1B_1 + \cdots + x_mB_m)\|^2 \\
    &= \|(c_1 - x_1)B_1 + \cdots + (c_m - x_m)B_m + R\|^2.
\end{align*}

Expandiendo la norma:
\begin{align*}
    \|A - x\|^2 &= \sum_{i=1}^m (c_i - x_i)^2 \|B_i\|^2 + \|R\|^2.
\end{align*}

Dado que $\|B_i\| = 1$ para todo $i$, esto se simplifica a:
\begin{align*}
    \|A - x\|^2 &= \sum_{i=1}^m (c_i - x_i)^2 + \|R\|^2.
\end{align*}

Para minimizar $\|A - x\|^2$, debemos minimizar $\sum_{i=1}^m (c_i - x_i)^2$. Esto ocurre cuando $x_i = c_i$ para todo $i$. En este caso, $\|A - x\|^2$ toma su valor mínimo, que es $\|R\|^2$. Por lo tanto:
\[
    \|A - (c_1B_1 + \cdots + c_mB_m)\| \leq \|A - (x_1B_1 + \cdots + x_mB_m)\|.
\]

Esto demuestra la desigualdad dada.